\documentclass[fontsize=12pt, appendixprefix=true]{scrreprt}
% appendixprefix: hogy odaírja, hogy "Függelék A", ne csak "A"
\usepackage[english, magyar]{babel}                        % nyelvi csomag
\usepackage[T1]{fontenc}                                   % ékezetes bet?knél is legyen automatikus elválasztás
\usepackage[utf8]{inputenc}                                % ékezetes bet?k kezelése
\usepackage{mathptmx,times}                                      % alapértelmezett bet?típus ne legyen pixeles
\usepackage{mathtools}                                     % képletekhez kell
\usepackage[style=ieee, backend=biber]{biblatex}           % bibliográfia
\addbibresource{thesis.bib}
\usepackage{graphicx}                                      % képek beszúrása
\usepackage[export]{adjustbox}                             % ez a logó pozicionálásához kell
\usepackage[margin=2.5cm, bindingoffset=0.5cm]{geometry}  % margók
\usepackage[onehalfspacing]{setspace}                      % másfeles sorköz
\usepackage[hidelinks, unicode, pdfusetitle]{hyperref}     % kattintható tartalomjegyzék és hivatkozások
\usepackage{bookmark}                                      % PDF könyvjelz?k
\usepackage{csquotes}                                      % a bibliográfiában megfelel?en legyenek formázva az idéz?jelek
\setcounter{tocdepth}{2}
\DeclareQuoteAlias{dutch}{magyar}

% Kódrészletekhez ajánlom
\usepackage{listings, scrhack}
\usepackage{sourcecodepro} % egy jó bet?típus
\lstset{captionpos=b, numberbychapter=false, basicstyle=\ttfamily, showstringspaces=false, columns=fullflexible}
\renewcommand\lstlistingname{kódrészlet}
\makeatletter
\renewcommand\fnum@lstlisting{\ifx\lst@@caption\@empty\else\thelstlisting.~\fi\lstlistingname}%
\makeatother

\titlehead{\includegraphics[width=1.0\linewidth,height=0.5\textheight]{logo}}
\subject{A természet megjelenítése a minta tervezésben}
\author{Sárkány Lili\\Kézműves Tárgykúltúra}
\title{Szakdolgozat}
\date{2021}
\publishers{Témavezető:\\Papp Anett}

% Nyilatkozathoz két parancs definíciója
\newcommand{\pushtobottom}{\vspace*{\fill}}
\newcommand{\signatureline}[1]{\begin{flushright}
	\vspace*{.5cm}\par\noindent\makebox[2.5in]{\hrulefill}
	\par\noindent\makebox[2.5in][c]{#1}
	\end{flushright}
}

\begin{document}
\maketitle

Alulírott ...
\pushtobottom
\signatureline{Aláírás}

\tableofcontents

\chapter{A természet mint eszköz és forrása}
Mint ahogy azt Kovács Nemere is kifejti a A képzőművészet története III. részében az ősművészeteknél \cite{nemerekepzHomHuveszet} a természeti motívumok és eszközök már az ősidők óta részét képezik az emberi művészetnek. Amikor a barlang rajzokat készítették számukra
az elsődleges forrás a természet volt.
Abban az időben az ember csak a környezetükben található tárgyakat tudták felhasználni és az akkor kialakuló emberi kreativitás által az egyszerű eszközök és lapanyagokból létrehozni a művészetet. A maguk tapasztalás útján jöttek rá, hogy  k miként tudják fel használni az őket körülvevő anyagokat: növények, ásványok, rovarok és mindent amit megismert az elméjük. Az így megismert és létrehozott alapanyagokkal kezdték újrateremteni az érzékelt világot primitív művészeti formákban, például barlang rajzok, kezdetleges bőr mívek vagy éppen primitív szövetek formályában. A legleső szín ami az ősi művészetekben megjelen a sárga és annak árnyalatai voltak amit különböző növényekből tudtak elő állítani.

\section{Alap fogalmak}
A dolgozat írás közben felhasznált alapvető és szükséges fogalmak melyek elengedhetetlenek a dolgozat egészében:
\begin{itemize}
	\item \textbf{Motívum} \\  A művészetekben a legkisebb önálló kifejezőegység, amely a mű során általában ismétlődik. A folklór - népművészet alkotásainak legkisebb tartalmi egysége, amely átvétel után is felismerhető marad.
	\item \textbf{Organikus} \\ szerves, növényi , állati eredetű alkotó elem, vagy művészeti motívum
	\item \textbf{Ornamentika} \\ (a lat. orno, 'felszerel, díszít' szóból): egy-egy kor, nép, műfaj vagy stílus díszítményeinek összessége. - Lehet mértani, növényi, figurális v. vegyes összetételű.
	\item \textbf{Festék} \\ Tágabb, köznapi értelemben mindenféle anyag, ami egy tárgy színét megváltoztatja: fedőfesték, máz, zománc, lakk, pasztell, lazúr stb. Szűkebb értelemben: a felületre felvitt, az alapot (szinte) teljesen elfedő, vékony, pigmentált réteg. A kötőanyagtól függően lehet: enyv-, olaj-, akril*-, tempera**-, viasz-, méz-, gyanta- stb. festék.
	\item \textbf{Színezék} \\
	Vízben oldható színezőanyag, mely kötőanyagot nem tartalmaz, jellemzően a kelmefestés kelléke. Az anyagot nem fedi el a színezék, hanem beleivódik a rostjaiba, a pigment kémiailag kötődik az anyaghoz. A létrejövő szín és teltsége függ a textilanyag fehérségétől, a kezelés időtartamától, a festőoldat sűrűségétől, az elő- és utókezelés módjától, anyagától (pácoktól*), ezek vegyi összetételétől, olykor más tényezőktől is.
	\item \textbf{Kelmefestés} \\
	\item \textbf{Kékfestés} \\ Egy textíliákon alkalmazott színmintázási technológia, amely nevét onnan kapta, hogy a minta eredeti formájában jellegzetesen kék alapon fehér színben jelenik meg.
	\end{itemize}

\section{A természet mint művészeti inspiráció}
A művészet utánozhatja a természetet azáltal, hogy egy absztrakción keresztül újra értelmezi azt, vagy képes teljesen lemásolni.
Művészet kinyit olyan ajtókat nekünk a természet felé ami a természet  bonyolultságra és szépségére hívja fel figyelmünket, amit semmi képen nem lehet elszalasztani.

Lehet egy egyszerű művészeti kép ami segít értelmezi a természetet vagy lehet egy kihívást jelentő darab ami kifejezi a emberi kapcsolat a természettel.
A természet által való művészeti gondolkodás lehetőséget add nekünk egy olyan látás mód kialakítására ami maga a természetben előforduló anyagokkal tárgyakkal való munkát helyezi a művészeti folyamat közép pontjába mint például levelek, botok ágak, termések, víz, kövek és ehhez hasonló természeti tárgyak használata.
Kreatív módon való felhasználást késztetést érzünk új művészeti tárgyak készítésére.

Ez a fajta művészeti filozófia fenntartja a kapcsolatot az őskori emberekkel és a természettel.

Manapság néhány kortás művész messze van attól a folyamtól ténylegesen, hogy saját mag hozzon létre festéket vagy bármilyen alapanyagot a saját munkájának a megvalósításában.

Az egyik közismert mai kortás alkotó aki művészetében vissza nyúl az ősi alapanyagokhoz Andy Goldsworthy. Olyan természetes anyagokat használt fel munkái során mint például levelek, kövek és hely specifikus szobrokat készített amelyek tükrözik az anyagok és a természet közöti őrsi kapcsolatot.

Munkáinak az elkészítési ideje hosszú folyamat, a felhasznált alap anyagok össze gyűjtése és előkészítése mélyíti el azt az ősi kapcsolatot ami művész és műalkotás között ki tudd alakulni.
A kész művei fényképformájában maradtak fent mivel ezek az alkotások a természetben helyezkednek el és semmilyen olyan alkotó elemet nem tartalmaztak ami megóvná a természet viszontagságaitól és lassan elamortizálódtak ilyen módon fejezte ki, hogy a tényleges műalkotás múlandó.


%Ezt majd át kell vinni
Dolgozatomban beszeretnék mutatni egy olyan kortás művészeti irányzatot, technológiát amelynek alapja a természetes anyagok felhasználásában gyökerezik és része kultúránknak és népművészetünknek.

% \section{A természet mint eszköz a művészetben}
% \subsection{Színanyagok}
% A természetben előforduló pigment anyagok melyek felhasználásra kerülnek a legkülönfélébb művészeti alkotásokban:
% Indigó, Bíbor, 
% \subsection{Eszközök}
% Nyomódúc. textilminták előállítására használt fa mintanyomó eszköz, leginkább puha és keményebb fa fajtákból készült mint például gesztenye fa vagy dió fa vagy tölgy vagy bükk


\chapter{Az indigóval való festés}
Az indigó név Római indicum -ból származik ami \textit{Indiai} terméket jelent.
Ez a megnevezés bizonyos értelemben téves név mert az indigót tartalmazó növényeket a történelem során a világ több pontján is termesztették, mint például Kína, Jáva, Japán és Közép-Amerika.  

Az 13. százelején amikor Marco Polo haza érkezet ázsiai utazásáról, amelynek során rájött, hogy az Europában hőn áhított és nagyon kereslet indigó nem egy ásvány hanem egy növényekből származó festékanyag és magával hozott egy palántát új színtörtbe az Európai divatba. Ezt követően már Európában is beszerezhetővé vált.

Majd a 16. században Portugália és Spanyolország bonyolította le az indigó kereskedelem nagy részét, mondhatni Indigó nagyhatalmakká lettek.
Ezt követően a Hollandok sajátították ki maguknak a mert elkezdték nagyobb mennyiségben szállítani az indigót köszönhetően a modernebb és fejlettebb hajó flottáiknak. Ennek következményeként tönkre ment az Európai  csüllengtermelés. A legjelentősebb mennyiség gyártást indiában kezdték el 1600 évek végén sok indigót hoztak be Európa.
 
\section{Technologiája}
 \textbf{Növényi indigó}
 Europában az indigó előállításánál festőfűként nevezet (festő csüllenget) használtak elsődlegesen. Ez a egy kétnyári fű fajta ami az év első felében levél szerű rozettákat nőveszt.
 Amikor az indigót kezdik el termeszteni, akkor learatják ezeket a rozetta szerűségeket. A következő évben virágkocsányokat és magokat hoz létre.
 A középkor tájékán fontos gazdasági szerepe volt a festő fűnek Europában ez hozta a legtöbb jövedelmet.  
 A hagyományos indigó festék előállításánál a festő fű leveleit péppé zúzták labda formákat készítettek belőlük ezeket több hétig szárították és ezt majd egy fermentációs, erjesztési folyamat követte majd a feldolgozás további lépéseiben is többször újra és újra össze zúzták.
 Sajnos a festő fű nagyon szennyezett volt és ezért csak egy nagyon halvány kék színt tudót adni.
 Ezzel ellentétben a trópusokról származó indigó sokkal jobb minőségű volt és mélyebb gazdagabb kék színt hozott.
 A növényi anyag kapott egy vizes áztatást utána elkezdték az erjesztési folyamatot és ezt követte a levegőn történő oxidációs folyamat. Maga a növény nem tartalmazta az indigót, hanem azokat az alap vegyületeket melyek a fermentáció során átalakulva hozzák létre a kívánt pigmentet.

\section{Földrajzi varianciája}
\subsection{Magyarországon}
A kékfestő mesterség fénykorábban 400 kékfestő volt Magyarországon ebből 300 a Dunántúlon volt ebből több fen maradó 100 b fen marado 60-70 kékfestö duna tiszta közén telepedt le és talan 5-10 jutott tiszán tulra a Tothok lakta vidékeken telepedtek  le Békéscsabán , Szarvasobn , Debrecenben 
10 több kék festö nem voltak város szerte 
nagyon sokan voltak akiki egyedül végztek el ezt a munkát mert nem volt elég helyuk hogy többen is elvégzék a feladatokat .
Mint például Pápán az egy sokkal nagyobb mühely volt és ott körübelül 30 dolgoztak 

A 18 század tájékán Mária Terézia idején kezdődött el Magyarországon az indigó és csülleng fű termesztés.
Azért próbálkozott Magyarország a termesztéssel mert szerették volna csökkenteni az Európát elárasztó indiai és türingiai festéket.
Ebben az időben komoly gazdasági alapja volt Magyarországon több mint 300 falu foglalkozott a csülleng fű termesztésével  és feldolgozásával.
   
% .......
% középkor idején elsősorban 

\subsection{Európában}
A kékfestett textil kereskedelem fő központja 18. században Franciaország  és Németország volt.

\textit{Franciaországban} a kékfestést technológiája eltért a többi Európai országhoz képest. Az ottani módszer jobban hasonlított a batikoláshoz. A textílanyag felületét viasz réteggel vonják be arra a részre ahol nem szeretnék hogy befogja a festék és ezzel a módszerrel sokkal tisztább és élesebb dekoratív mintákat tudtak elérni ennek a viaszos alkalmazásnak a segítségével.
Mintakincsük sokkal aprólékosabb részlet gazdagabb inkább levélszerű mintákat alkalmaznak, ezek közt elvétve fel tűnik egy két virág motívum mely apró és elhanyagolható a teljes mintában.
A minta fa készítéshez gyümölcs fa fajtákat használtak fel mert sokkal lágyabbnak és puhábbnak tartják a véséshez és kaparáshoz. Az általuk előállított anyagok több színűek és esetenként a minták is színesek.
Például a vörösszín eléréshez festő fűt alkalmaztak a festő buzér - Garance - Rubia tinctorum  ez a nővény alapozta meg a vörös szín használatát, 
a sárga és a zöld szín előállítása a festő csülleng feldolgozásának részfolyamatainál mint melléktermék jött elő.

\textit{Németország} 
Itt három technológia volt elterjedve a Tartalék (Reservedruck), közvetlen (Direktdruck), maratás (Ätzdruck). Az alap texílt először fehérítették, szárították és végül hengerelve (Mangeln) vasalták és keményítették.

A tartalék nyomás során fehér mintáthoznak létre kék alapon. Az eljárás során a gumiból hozzák létre a nyomatot. Eztkövetően a nyomatot (Reservage) bevonják egy festék tarszító anyaggal és a minta nyomás száradást követően történhet. A festék öszetétele: fehér dohány, réz-szulfát, réz-acetát és egéb titkos összetevők. Pasztell vagy német indigónak hívták ezt a festéket. A festés után a mintát higítot kénsavval távolítják el.

Ezért ez egy festési eljárás, nem egy nyomtatási folyamat. A "tartalék" kifejezés arra a tényre utal, hogy a kiválasztott minta a festés során kimarad.

A közvetlen nyomtatásban a színt az előkészített fehér szövetre nyomtatják a nyomdarúddal. A szín közvetlenül az anyag felületére kerül, és barna színként jelenik meg. A száradás után a szövet egy fürdőbe helyezik, amelyben a barna tinta kémiai reakcióval élénkkékre változik. A szövet végül felfőzik, újra hengerelik, eztkövetően használatra kész. A nyomtatást nagyon óvatosan kell elvégezni, mivel a hibákat nem lehet kijavítani.

A maratás folyamata egy kék színű alap szövet is nyomtatható maró anyaggal (maró folt), ami kék alapon fehér mintához létre. Ez a technika alapvetően hollandiából származik, de máshol is elterjedt az eljárás.


\subsection{Ázsiában}
A \textit{Shibori} egy a 8. századból származó japán textil festési technika amelynek lényege, hogy különböző kötözési, csomózási eljárásokkal állítanak elő tetszőleges mintázatott.

A sibori számtalan más japán alap technikát foglal magában, iylen technikák például a miura, a kumo, a nui, mindegyikben közös, hogy valamilyen kötözési módszer vagy az anyag hajtogatást alkalmazz. 

A kötözés és halytogatás mellet a szövetett olykor festéktaszító anyaggal is bevonják és ezzel is mintákat hoznak létre.

A teljes minta elő állításának folyamata az alap szövet tisztításával és esetleg fehérítésével kezdődik, eztkövetően történik a vasalás és halytogatás vagy a szövet össze fogása és kötözése és felületének kezelése.
Az siboriban az hogy milyen technikát használunk nem csak a saját preferenciánktól függ vagy a vágyott mintától, hanem a befestendő anyag tulajdonságaitól is. Ezen kívül  a különböző festési módszerek kombinálhatunk kidolgozottabb, bonyolultabb, részletesebb minták elérése érdekében. 

\subsection{Különböző régiók összehasonlítása}
   Európa nyomódúc <> Batikolás

\chapter{A kékfestés}
\section{Előzménye és Történelmi áttekintése}
Kékfestészet előzményének a kelmefestészetet, és a textilnyomást tartjuk ami egyszínűre színezést és mintázást jelent mely technológiák már a 15. - 16. században is léteztek. Sok tudást igénylő művészeti ág volt.

A középkori kelmefestés kezdetekben a városokban és a kolostorokban működtek különféle anyagok színezésével foglalkoztak mint például gyapjú vagy vászon. 
Az indigó megjelenéséig magyarországon a kelmefestésre a festő csülleng növényből elő állítot kék festékanyagot használták. Ez a 15. -ban kerül magyarországra és terjedt el fő alapanyagként.

A csülleng fontos szerepet töltőt be a középkortól kezdve, legnagyobb termelő területek Franciaország és Türningában volt. A magyarországi mestereknek a beszerzési központ Türnigiában volt.
A 16 század körül körülbelül 300 falu foglalkozott a csülleng termesztéssel Magyarországon. Aki akkoriban a csülleng termesztéssel foglakozott vagyonosabb réteghez tartozott. A 17 századra vissza eset a termesztés mert ekkortájt a holland hajókaravánok behozták kelet indíából az indigót.
Ezekben az időkben már csak kűrű belül 30 falú foglakozott még a termesztéssel.

\textbf{nézd át}

Brassó város könyveiben jelent meg 1600-as  évek körül  az első feljegyzés   a posztó kereskedelemről .
Létezett  egy céhen kívül szervezet aki szét osztotta a mesterek között festendő  posztót és az indigót .
 miután a mesterek el készültek a festetett mintás anyagokkal azokat összeszedték és ki vitték őket vársárokra 
Az országnak azon északi részén ahol még nem voltak törökök 16-17 század környékén hirtelen fejlődés vette kezdetét 
De akik menekültek a törökök elöl más falvakba kényszerültek le telepedni , és ezek a falvak már kicsit zsúfolttá váltak .
a termelés és maga a munka folyamatok túl zsúfolttá váltak    .
A kisvárosokban mint például Lőcse , Eperjes itt már a 16 században több mester  fogalakozott anyagok és fonalak  festésével  .
AZ 1608-as évek idején sok vándorló volt és a magyarországi városokban más ott élő mesterekkel közösen hoztak létre festő újabb céheket .

Az első magyarországi feljegyzés kékfestésről 1783-bol  van Körmöcbányáról 

17 század közepe táján  meg is honosodott kékfestészet Magyarországon 

18 században szinte minden városban volt olyan szak ember aki ezzel foglalkozott mert igen nagy kereslet vette a kezdettét nyomot mintás anyagok iránt sok mester próbált új technikákat létre hozni , illetve  elsajátítani amik már léteztek .
 Sopronban volt egy  család akik  19 században ,  fontos szerepet töltöttek be  mert ők voltak az első mesterek akik gyár jellegű kékfestő üzemet működtetek ,
 Aztán 1852 ben és 1861 ben  elkezdték  a bővítést , és a későbbiekben külföldre exportálták termékeiket  Európa minden  pontjára .és  Indiában béreltek földeket ahol saját maguk termesztették indigót  .
A  Goldberger, és a Szegedi  műhely is , külföldi tapasztalatok alapján végezte  ezeket a  munkafolyamatokat .
  Aki a Szegedi műhelyt vezet az  Felmayer Antal volt aki 1826-ra gyárrá alakította , fejlesztette a műhelyét.
Amikor megtörtént az ipari fejlődés egyszer csak maga a manufaktúra és az ipar kezdet úgy mond a végéhez közeledni de próbáták még benne tartani az erőt .
A vidéki életformák változása után a városi dolgozóknak az öltözködésük is meg változott a régi hagyományos viseletek igényei is változni kezdtek Magyarországon 
Az első és a másodig világháború után nyersanyag hiány lépet fel városiasodás csökkentette a műhelyek számát 
De még egy két vidéki helyen ma is foglalkoznak és viszik tovább ezt a mesterséget próbálják fent tartani ezt hogy ne veszen el a kékfestészetnek maga a varázsa .

\textbf{idáig}

\section{Technológiája}
Egy kékfestő műhely létrehozásához is jelentős tőkebefektetésre volt szűkség.
  Hiszen egy teljesen külön épület kellet hozzá .
Volt egy helység ami a színezés és a kifőzés helyet szolgálta , és  fekete konyhában folytak a kifőzési folyamatok.
Minta készítés helye a küpa szobában történt.
száritás az udvaron folyt ha rosszabb  idő volt ezt akkor bent kellet valahogy kivitelezni .
 Ezek a műhelyek általában   a család telkekre épültek  . 
 A A nagyobb műhelyek rendelkeztek körű belül 4-5 katlannal .
 A vásárolt vásznakat  ki főzték a a szennyeződések miatt , másfél órán keresztül szódás vízben főzték .
 Miközben ment a kifőzés ,  keményítős oldalát meg jelölték nehogy arra az oldalára kezdjék el a mintázást, mert sajnos attól hogy kifőzték előfordult hogy foltosan fogta be az anyagot ,
 A fen maradt végeket fekete konyhába vitték  körű belül , egy literes fa kádban tiszta vízben ki öblítették .
 Miután befejeződött a öblítés jött a száritás.
 Ezt a folyamatot udvaron történt meg , száritó az úgy nézet ki , hogy  az udvar közepén ált lábakon  vagy az épület oldalára volt erősítve . De mikor meg érkezett a fagyos időszak akkor sajnos nem volt lehetséges a kinti száritásmert nem volt jó anyagok számár mert  a keményitő  ,elkezdet el színeződni az anyagon  .
 Leginkább ez az esett jobban a festet , anyagokon látszódott meg sajnos és azt már nem tudták felhasználni ez selejtek közé került ,
  ebben az időszakban szüneteltetniük  kellet a munkát körű belül 6 hónap kimaradás volt a részűkről
Igy akkor kevesebb munkásra volt szűkségük de addig az úgy mond pihenő idő  alatt se tétlenkedtek hanem a minta dúcokat tudták előkészíteni addig is .
 A századforduló idejére már több vidéki műhely is rendelkezett gőzgéppel vagy kazánházzal ezeknek a segítségével már betudtak szerelni télire szárítókat  a műhelyen belülre is.
   Búzát , kukoricát  , burgonyát alkalmaztak az anyagók keményítésére , ez a folyamat a fekete konyhában zajlott le .
 A megszáradt végeket elvitték a Mángorlóba , ez a  szerkezet  egy erős gerendára szerelt asztalalapzatból állt  rajta erős vastag keményfa, görgőkön mozgót ,mérete körű belül 4-6 méter hosszú lehete és körű belül két és félméter magas lehetet
 akkoriban ezeket lovakkal hajtották .
 Kapcsolódott hozzá egy láda ,  a fölött  el haladt egy főtengely és amire  került fel egy vastag lánc ami az asztalon lévő ládát mozgatta .
 A belső kerekeknél két főtengely helyezkedett el váltókar segítségével lehetett irányitani ezzel azt tudták elérni hogy jobbra és balra is tudták forgatni a szerkezetet.
 Plusz  az volt benne hogy kis helyen is elvéghető legyen ez a munka folyamat  .
 Előfordult olyan is amikor viSsza esett egy kicsit a az Üzlet forgalma  és akkor sajnos nem volt pénz lóra aki hajtsa a kereket akkor maga a mester álltbe és tekerte a szerkezetet
 A mángorló több változta is létezik az általános méret 30x30  ezeket a szerkezetek általában vastag gerendákból építették fel 15-20mm vaslemezeket szereltek fel rá .
 Azokba a műhelyekben ahol adott volt a gépesítés ott gőzgép vagy transzmissziós tengelyekkel vitte az erőt a mángorlóba ezt a technológiát helyileg ,Pápán ,Csornán, és  Békéscsabán használták .
  Más varosokban ( Bátaszék) előfordult  hogy itt benzinmotort  alkalmaztak Csornán villanymeghajtásura építettek át a mángorlót ,
 Fontos alkatrésze a mángorlásnak a görgő , és a áruhenger volt .
 Az únt felhelyezték a felsodrószékre  és a vászon végét beburkolták .
 Miközben mozgásban volt nagy súly volt rajta , és ide oda mozgott  eközben ki simította az anyagot .
 Mielőtt elkezdődött a mintázás kétszer háromszor végezték el ezt a folyamatot festés után a keményítés körű belül hat - nyolc-tíz alkalommal végezték el 
 Maga a mángorló téglával lerakott helység volt itt készítették el a vegyszereket is .
 Amit a minta készítésnél használt fő fedő anyag az a pap volt álltlában más - más összetételben készült ez anyag nagy titkos receptek alapján készítették el 
 és hogy több hónapig elegendő legyen igy egyszere egy nagy mennyiségek készítettek el 
 Ahhoz hogy ne mennyen tönkre ez az anyag nyirkos hűvös helyen tartották 
 ezt az anyagot különböző  pácokhoz és fedőanyagokhoz is hozzá keverték mikor megkezdődött a kék festés akkor mész lúgos öblítéssel sárga zöld , narancs színű mintákat is tudtak előállítani .
 Nálunk nem volt megszokott színes mintás anyagok készítése mi maradtunk fehér minták létrehozásánál ez a stílus inkább a franciák körében volt  jellemzőbb.
 A mintázáshoz is volt egy külön helység amit használtak és ott abban a helységben volt a tarkázó asztal , e mellet helyezkedett  el a a Sasi.
 A Sasi egy olyan láda volt ami körülbelül 6-8 cm mélységű volt , és benne  egy sűrű keményitő .
 A ládában volt egy keret aminek ez egyik oldala viaszos vászonnal be volt vonva  a másik  része pedig monílóval , van be vonva ebben kenik el a fedő masszát a papot  és kürű belül úgy működik mint a pecsét párna .
 A minta készítés nyomódúcokkal folyt .
  Méret különbségek voltak a  dúcok között, de ahhoz hogy hatékonyabb legyen  mintás anyagok készítése nagyobb dúcokra kellet válltani de azoknak az elkészítési ideje sok időt vett  igénybe igy át váltottak egy  úgy nevezett perrotingépre ami fel gyorsította a gyártást .
 1860 körül járhatunk amikor  ugyan is bekövetkezett a századforduló és a gépiesítés igy egyre több műhely állt át  a gépi nyomásra .




\subsection{Alapanyaga}




\subsection{Folyamata}
kékfestészet  munkafolyamata úgy zajlott le egy műhelyben, hogy a nyersanyagot magánszemélytől vásárolták ,  vagy maga a megrendelő vitte magával az anyagot és még a mintát is ki tudta választani  a minta könyvből
Budapesten voltak olyan emberek hogy csak azzal foglalkoztak a kékfestéssel foglakozó mestereket látták el anyagokkal.
Az első  folyamat amit kezdenek az a anyagokkal, hogy ki mossák őket szódás vízben ez körülbelül két órát vesz igénybe.
 A szódás kifőzést azért végzik el hogy az anyagból ki oldódjon egy gyantás anyag és ez a folyamat után  az anyag kifehéredik .
 Miután ez a folyamat befejeződött akkor következik a hideg vizes öblítés majd léces szárítóra helyezik majd néhány óra alatt meg szárad 
Mikor már megszáradt az anyag  kap egy vékony keményítést ami előkészíti a mintázáshoz 
Mintázást régen úgy nevezték  hogy Tarkázás.

Tarkázásnak említett kézi minta készítés  párnázott asztalon készül.
Amire több rétegben pokrócokat helyeztek majd a végén molinóval befedték igy az asztal felületét úgy lehetet elképzelni mint a pecsét párna lenne .
Igy az asztal felülete nagyon puha lett.

Az asztal mellet helyezték el Únt és az ún-ba mártották bele a minta dúcokat.
A viaszosvászon belső felületére kent fel a mintaázó pépet az Ún 
ládát feltöltötték sűrű keményitővel ami a sasit puhán tartotta .
Amit általában használunk hozzá  fehér nyomó pépet annak alkotórészei ólom nitrát, ólomacetát . gumiarábikum . kék.timsó, rézgálic .
Ezeket össze főzik és miután ki hűlt egy hétre rá használják csak fel  az anyagot 
A festék pépet  maga a mester szokta elkészíteni általában ez a festék pép akár évekig is elég szokott lenni .
 Mesterek  az gondolják hogy egy pép minél régebbi egy  annál fehérebbé teszi a mintát
A kékfestő a minta  dúcot a nyomópépbe nyomja majd és a dúcot vászonra illeszti pontosan nehogy elcsúszon és majd erősen lenyomja hogy minden pontra kerüljön a pépből és azt a részt nem fogja be a kék  indigófesték .
Mintázás nagy figyelmet és türelmet igényelt.

Általában a kékfestők  ,120-130 méter anyagot meg mintáztak , de az napi 13-14 órás munka folyamat igényelt. 
Miután a minta nyomás befejeződött  párnap volt még megszáradt a nyomatás ezek után következet a festés.
Az anyagot egy ráfra akasztják ,  ami  igazából egy vaskerék érre akasztják fel az anyagot .
Ez tehát  úgy néz ki  ez az anyag a ráfon  olyan az esése  mintha vaskereken egy rakott szoknya lenne .
Következő teendő pedig nem  más volt hogy elkezdték  lengedik  az indigó festékkel teli kádat és  félóránkként fel és le mozgatják az anyagot  , és körülbelül - percet hagyják a levegőn hogy az indigó oxidálódjon .
Amikor elsőnek felhúzzák az anyagot a festékből akkor ilyen sárgás zöld színe van aztán megint eltelt egy kis idő akkor ilyen sötétzöld színe lesz az útólsó előtti merítés után világos kék színe lesz 
miután egyre többet éri a levegő  megkapja azt a szép sötétkék színt 
Miután az anyag elérte a kívánt színt akkor következik a száritás. Ezután miután meg száradt az anyag akkor kap egy kénsavas fürdőt ennél a folyamatnál az történik hogy kioldódióik a mintázásnál bele ragadt  nyomópép és ezek után nagyon szépen elő jön a minta .
 Utána egy sima vízben is ki öblítik  és hagyják meg száradni , ki keményítik az anyagot .
Ez volt a teljes munkafolyamat.


\section{A természet mint inspiráció}


\section{Minta dúcok}
Maga a minta és a forma készítés rajztudást igényelt és a textilnyomáshoz való ismereteket igényelt 
mintáknak mindig pontosan kellet egymáshoz illeszkedniük  ezt a készítőnek mindig tudnia kellet .
A nyomódúcok felületei általában mindig kiemelkedtek és negatívba kellet rajtuk a mintákat elkészíteni 
Két fajta minta dúcot készítünk fából kifaragott vagy réz lemez formájú mintát 
első az hogy  fából kifaragott mintákat ki rajzolják a mintafára utána azt elkezdik kifaragni 
a következő fajta úgy készült hogy a fatáblát kivésték  rézlemezek helyét fémlemezeket a helyükre tették és utána vízbe áztatták 
A dúcok hátlapjára egy erős fa lemezt erősítettek és ez a fa lemez segítségével könnybe volt a nyomdázás
A kedvenc és bevált mintákat idomait mintázó gép nyomó fáira is átvitték .
Igy ma készített minták ornamentikája őrzik a több évtizedes mintákat és ízlés hagyományokat.
A kezdetekor úgy nevezték a kézi minta rajzokat mint például  ágas, ágasindás , ágasvirágos, csokros, cakkos ,folyondár , margarétás , négyzetpontos ,
akkoriban nagyon egyszerű neveket adtak a kész mintáknak .

A 17 század és a 18 századi minta dúcokon pozitív  nyomást alkalmaztak amikor növényi vagy figurális mintákat készítettek ezeket a fa metszeték  dúcokra készítették.
Ugyan ebben az időszakban már szinte minden városban volt egy minta készítő mester .
Sasvári a Cseklésziek  a 18 század közepe táján ők több minta készítő mestert alkalmaztak azért volt annyi emberük hogy minél gyorsabban el tudják készíteni a mintáikat 
még az itteni munka mellet tudtak foglalkozni vidéki műhelyek mintáival de miután végeztek mentek is tovább .
Miután elkezdtek terjedni nyomot kelmék igy vidéki műhelyek mintakincsei is kezdtek felkapotabbá  válni .
Ez idő alatt külföldről be érkező emberek letelepedtek egy adott időre az adott mestereknél akinek elkészítette a mintáját és addig kapott ellátást neki csak a munkával kellet foglalkozni miután végzett ment is tovább 
Ez idő alatt Magyarországon csak ez a kettő minta készítő műhely volt már korábban említett Sasvári és Cseklészi műhelyek vagy úgy nevezett manufaktúrák voltak .
Későbbiek ben jöttek létre nagyobb gyárak mint például Goldbergerés Spitzer Gerson és aSzegedi Felmayer ők alkalmaztak mintafaragó és készítő szakembereket 
legelső minta fák teljes mértékben fából készültek kisebb pontok, ágak , hasonló képen mint a folt felületek de ezeknek a készítésére legalkalmasabbnak a dió fát tartották vagy körte és gesztenye fa mert ezek a fajták nem szálkásodtak  egyszerű volt ki faragani belőlük a mintákat tartósak voltak és  nem repedeztek ki .
A 16  század idején néhány nyomó dúc kőzött elfordult tölgy és  bükk is ezeknek használatával sok változás jött elő
ugyan is ezeknél kezdték el használni a réz drótot de későbbiekben már teljes mértékben vörös réz drótot használtak
azért volt jobb csak a vörös réz drótot használni mert az előállítása olcsóbb volt és egyre több mintát tudták elkészíteni 
A fémből készített mintáknak a legfontosabb tartozéka az ún volt .
A fémből készült minták nagyon időigényesek voltak 
A mintákat úgy kezdik el megtervezni hogy elszőr papír terveket készítenek miután a elfogadták a mintát akkor elkezdi felrajzolni a minta fára és utána be lakozza hogy miközben készít ne mosódjon el rajt a minta .
ritkább vagy bonyolultabb mintáknál a be karcolják a a mintafát hogy pontosan helyezkedjenek el az elemek .
A minta dúcokra készített mintákhoz lapos vésőt ütögetve haladnak  és közben beverik a a lemezeket .Ha több szint szeretek volna alkalmazni akkor egy kiegészítő dúcot alkalmazták 
A békéscsabai festő műhely tulajdonosa Sztaricskay Pál  régi kedven kézi mintáit elkészítette a gépi nyomáshoz is .
Még aki gépi mintákat készített az Gál Gyula volt és ö még adott is el belőlük .

\subsection{Alap minták}
Az öltözködés terén többnyire női ruhákon lehetett látni könnyed és nagyobb ,egyszerűbb  alap mintákat formákat mint például virágók , levelek , apró pöttyök vagy geometrikus formák .
de ezek a minták ugyan úgy rá kerültek ágymenükre zsebkendőre de terítőn is felehetők voltak ezek könnyed minták .
A természeti motívumok mellet kezdtek elő jönni a geometrikus formák is.
A kezdetekor sokkal egyszerűbb mintafákat készítettek mert akkor még nem láttak és még nem is tudták mire lesz majd igény és kereslet az évek előre haladásával.

\subsection{Minta fejlődés}
Minták olyan irányban fejlődtek tovább hogy nem csak már említett egyszerűbb minták voltak hanem már azokat is tovább vitték.
 Sokkal nagyon mintafákat készítettek sokkal sűrűbb és aprólékosabb virág, levél  motivumokal készítettek voltak külön megrendelések a későbbiekben már voltak olyanok akik családi címerüket szerették  volna minta fára
 készítettek olyan mintát ahol már történtek szerepeltek mint például Ádám és Éva története vagy vadász jelenetek voltak meg jelenítve az nem minden esetben volt tiszta hogy igazából miért is jelenítetek meg eseményeket
 Pápán1 16 század közepe tájékán  készítettek hímzést untánzó kendő mintákat is .
  Többi város városban  mint például Gyula itt damasztszövéshez hasonló minta fát csináltak és még  , Cegléd  az ottani mesterek is  készítettek szövést és hímzést ábrázoló mintákat .


 


\subsection{Egyébb fő minták motívumok}
\section{Hagyományos kékfestő műhelyek}
Tolnai Kékfestő műhely Magyarország  egyik  legrégebbi még ma is működő kékfestő műhely 1810-ben alpítoták  .
Műhelyben még fellelhetők a régi eszközök  az eredeti hagyományokat követik még a mai napig a régi eszközökkel  dolgoznak .
Motívumokra jellemző hogy csak is régi családi mintafákat használják fel .
Ezekből a minta készletből több mint ötszáz darabbal rendelkezik a család .
Sajnos ezek a darabok nem pótólhatok mert hiszen hazákban évtizedek óta kihalt ez a mesterség , nincsenek már olyan mesterek akíkk  a minta készítéssel foglalkoznak 
igy nincs nagyon lehetőség pótolni a tönkre ment darabokat  
\section{Egyéb műhelyek}
\subsection{Kortás megoldások}
\subsection{Adaptációk}

\chapter{Összefoglalás} 


\printbibliography
\renewcommand{\appendixname}{Függelék}  % különben "Appendix"-et ír
\appendix
\chapter{Mellékletek listája}

\end{document}